\section{C. Hole filling, quotient actions, and exterior groups}\label{part:C}

We retain the notation $r_s:H_s\to G$, $q_s:H_s\twoheadrightarrow Q_s$
of Part B. The word ``pushout'' below means a bicategorical
pushout: cones include invertible comparison two-cells, and
the universal property is an equivalence of Hom categories.
The mixed-topos triangle has a different, noninvertible
comparison cell. This distinction is explicit on LM
scans~164--166.

\subsection{The two squares and the two kinds of gluing}

\paragraph{C.1 (Proved): the squares of (**) and (****).}\label{C1}
\textit{Manuscript: LM \S16, folio 81--83, scans 163--168.}
Let $\mathcal B_D$ be the gluing of $\coprod_s\beta H_s$ and
$\coprod_s\beta Q_s$ along $\prod_s q_{s*}$, and let
$R:\mathcal B_D\to\mathcal B_{X;S}$ be induced by the $r_s$.
Both squares in
\[
\begin{tikzcd}[column sep=2.6em]
 \coprod_s\beta H_s\ar[r]\ar[d]&\mathcal B_D\ar[d,"R"]
                      &\coprod_s\beta Q_s\ar[l]\ar[d,equal]\\
 \beta G\ar[r]&\mathcal B_{X;S}&\coprod_s\beta Q_s\ar[l]
\end{tikzcd}
\]
are two-cartesian in the two-category of all Grothendieck
topoi. The same is true for the open and closed squares in
the geometric diagram (****). For a henselian trait its
entire etale topos is the indicated $\mathcal B_D$.
The upper triangle through $q_s$ has only its canonical
\emph{lax} cell; it is generally not isocommutative.

\begin{proof}
In triples, $R^*(M,N,a)=(r^*M,N,a)$. The open
subterminal $(1,0)$ pulls back to $(1,0)$. A geometric
morphism to a topos factors through its open subtopos
precisely when that subterminal pulls back to $1$, and
through the complementary closed subtopos precisely
when it pulls back to $0$. These two universal
properties identify both pullbacks and their comparison
two-cells. They apply equally to the inverse image of
the open $U$ and closed $S$ in $D$.

Every sheaf on a trait is determined by its restriction
to its generic field, its restriction to its closed
field, and the specialization map from the latter to
the inertia invariants of the former. This follows by
the ordinary open--closed sheaf gluing construction:
on an etale object take pairs of sections whose germs
agree on the generic parts of its local branches.
The sheaf condition makes these pairs glue uniquely,
and restriction recovers the original pair. The generic
and closed field topoi are $\beta H_s$ and $\beta Q_s$,
and the gluing functor is $q_{s*}$. Thus it gives the
claimed identification. On a triple, the comparison
cell is exactly $a_s:q_s^*N_s\to r_s^*M$. It need not be
bijective: $N_s=\varnothing$, $M=1$ already disproves
invertibility. This verifies the warning of scan~164
and the two-square assertion of scan~166 separately.
\end{proof}

\paragraph{C.2 (Proved): the exact hole-filling pushout.}\label{C2}
\textit{Manuscript: LM \S16, folio 82--84, scans 166--169.}
For discrete groups put
\[
 \overline G=G/\langle\!\langle r_s(\ker q_s):s\in S\rangle\!\rangle.
                                                               \tag{C.2.1}
\]
For profinite groups take the \emph{closed} normal closure.
Then $\beta\overline G$ is the bicategorical pushout of
\[
 \beta G\longleftarrow\coprod_s\beta H_s
                    \longrightarrow\coprod_s\beta Q_s          \tag{C.2.2}
\]
in all ordinary Grothendieck topoi. Hence it is also the
pushout among locally connected topoi and among the
corresponding multi-Galois topoi. In the profinite case
the surjectivity of each $q_s$ is material to this
strong statement. Formula (C.2.1) is scan~169's formula,
and B.2 identifies it with the geometric filled group.

\begin{proof}
For a discrete groupoid $C$ and any topos $\mathcal T$,
\[
 \operatorname{Hom}_{\mathrm{Top}}(\beta C,\mathcal T)
                \simeq\operatorname{Fun}(C,\operatorname{Pt}(\mathcal T)).
                                                               \tag{C.2.3}
\]
Indeed the components of an inverse image
$\mathcal T\to\Fun(C,\Set)$ are point functors; the
$C$-actions supply their natural isomorphisms.
Conversely such a functor assembles a left-exact
cocontinuous inverse image. This describes also all
two-cells, including noninvertible natural
transformations between points. The pushout groupoid
of (C.2.2) is obtained by freely adjoining the
comparison isomorphisms and imposing the group
relations. Its functors into any category, and their
natural transformations, are precisely the
isocommutative cones. Since each $q_s$ is surjective,
one can eliminate all the $Q_s$ generators. The
remaining relations are $r_s(\ker q_s)=1$. The
result is the one-object groupoid $\overline G$.
Equation (C.2.3) proves the full topos universal
property, not merely its restriction to Galois targets.

Here is a direct proof which also treats profinite
groups. A cone into $\mathcal T$ consists of inverse
image functors into continuous $G$-sets and continuous
$Q_s$-sets, identified as $H_s$-set-valued functors.
After using that identification on the underlying
point functors, the $G$-action on every value kills
$r_s(\ker q_s)$. The kernel of an action on a discrete
set is closed, being the intersection of its open
point stabilizers. It therefore kills their closed
normal closure. The action factors continuously
through $\overline G$. Conversely an action of
$\overline G$ restricts along $H_s$ through $Q_s$;
surjectivity and the quotient topology of $q_s$
give precisely the required continuous $Q_s$-action.
The same argument on each natural transformation
proves full faithfulness of the universal equivalence
of Hom categories. Left exactness and preservation
of colimits are unchanged by the factorization.
This proves (C.2.2) in the stated largest two-category.
\end{proof}

\paragraph{C.3 (Refuted, with replacement): arbitrary profinite pushouts and higher shapes.}\label{C3}
\textit{Manuscript: LM \S16, folio 82--84, scans 166--169.}
Arbitrary pushouts of profinite groupoids do not
become pushouts in all ordinary topoi under
$G\mapsto\beta_{\mathrm{cont}}G$. Even locally connected
target topoi detect the failure. The correct general
statement is the pushout in the two-category of
profinite Galois topoi and continuous morphisms.
For the surjective right legs of the manuscript,
C.2 supplies the stronger result. None of these
ordinary two-pushouts is in general the homotopy
pushout of the full shapes.

\begin{proof}
Take the span $C_2\leftarrow1\to C_2$. Its profinite
pushout is $P=C_2\widehat * C_2$. Its discrete pushout
is the infinite dihedral group
$D_\infty=C_2*C_2=\langle a,b\mid a^2=b^2=1\rangle$.
The two inclusions give a cone of topoi into
$\beta D_\infty$. If $\beta_{\mathrm{cont}}P$ had
the asserted ordinary universal property, this
cone would induce a geometric morphism to
$\beta D_\infty$. Pull the regular $D_\infty$-torsor
back and choose a point of its underlying torsor.
The resulting action is a continuous homomorphism
$P\to D_\infty$ with discrete codomain. Its image is
finite by compactness, but the cone requires its
image to contain $a$ and $b$, which generate the
infinite group $D_\infty$. Contradiction. The target
$\beta D_\infty$ is locally connected. The normal
form $(ab)^n$ has length $2n$ for every positive $n$;
the supplied program checks the first ten cases,
while this reduced-word argument proves infinitude.
The example does not have a surjective right leg
and therefore does not contradict C.2.

Within profinite Galois topoi, continuous functors
of their fundamental profinite groupoids give
exactly the corresponding universal property;
the free profinite product is characterized by
continuous homomorphisms to profinite groups.
Finally, in discrete topoi
$\Set\amalg_{\beta\mathbb Z}\Set\simeq\Set$ by C.2,
whereas the homotopy pushout of their classifying
spaces is $*\amalg^h_{S^1}*\simeq S^2$.
Its second cohomology is nonzero. Thus a
two-pushout of covering topoi cannot be substituted
for a homotopy pushout of full spaces or shapes.
The finite simplicial calculation of the latter
cohomology is included in the program.
\end{proof}

\paragraph{C.4 (Proved): transitivity of partial filling.}\label{C4}
\textit{Manuscript: LM \S16, folio 84, scans 169--170.}
If $S=S_0\amalg S_1$, fill first $S_0$ and then the
images of the loops at $S_1$. The resulting group
and topos are canonically those obtained by
filling $S$ at once. The notation on scan~169
must read $U\cup S_0=X\setminus S_1$.

\begin{proof}
Let $N_i$ be the normal closure of the loops in
$S_i$, closed in the profinite case. The iterated
quotient is
\[
 (G/N_0)/\overline{\langle\!\langle\im(N_1)\rangle\!\rangle}
 \simeq G/\overline{\langle\!\langle N_0,N_1\rangle\!\rangle}.
\]
Both quotients represent exactly the homomorphisms
from $G$ killing both families. The same universal
description includes conjugators and hence gives
the canonical bicategorical equivalence, with
its commutation cells, for arbitrary groupoids.
In geometry $S_0$ has been added to $U$, so the
remaining punctures are $S_1$. This proves both
transitivity on scan~170 and the indicated
correction of its neighboring subscript.
\end{proof}

\subsection{Acyclicity and equivariant quotients}

\paragraph{C.5 (Proved): exact acyclicity statements.}\label{C5}
\textit{Manuscript: LM \S16, folio 83, scans 167.}
The map $\mathcal B_D\to\coprod_s\beta Q_s$ is
cohomologically acyclic for all abelian
coefficients pulled back from the target; in
fact $R\Gamma(\mathcal B_D,F)=\prod_sR\Gamma(Q_s,N_s)$
for an arbitrary abelian triple $F=(M_s,N_s,a_s)$.
The map $\mathcal B_{X;S}\to\beta\pi_1(X)$ need
not be acyclic. For a smooth geometrically
connected curve and torsion locally constant
coefficients from the target, it is acyclic
if $X$ is affine or has positive proper genus.
For proper geometric genus zero the complete
list, with $r=|S_{\bar k}|$, is:
\[
\begin{array}{c|c}
 \text{case}&\text{acyclicity for every torsion local system}\\\hline
 r=0&\text{yes, in every characteristic}\\
 \operatorname{char}k=0,\ r=1&\text{yes}\\
 \operatorname{char}k=0,\ r\geq2&\text{no; first geometric failure }H^2\\
 \operatorname{char}k=p>0,\ r\geq1
       &\text{yes for }p\text{-primary; no for prime-to-}p.
\end{array}                                                     \tag{C.5.1}
\]
In the failing cases the exact arithmetic obstruction
for a coefficient module $M$ is
\[
 R\Gamma(k,M^{(p')}(-1))[-2],
\]
with $M^{(p')}=M$
in characteristic zero. Thus vanishing of that
complex is the exact coefficientwise arithmetic
criterion. Acyclicity of the mixed-to-filled map
does not imply that the curve-to-mixed map is
acyclic.

\begin{proof}
For a local triple the global section is determined
by its $Q_s$-invariant closed section: its generic
section must be its image under $a_s$. The
closed evaluation has the exact left adjoint
$N\mapsto(q_s^*N,N,\id)$ on abelian objects.
It therefore preserves injectives. Deriving
the global section formula proves the first
assertion. Equivalently the first and third
terms of B.3's fibre diagram are identified
by the identity of $H_s$ and cancel.

When $X$ is affine or proper of positive genus,
B.5 computes its local-system cohomology by
$\pi_1(X)$, while B.6--B.7 compare it with
the mixed topos. This proves the positive
assertion. If $r=0$ the mixed and filled topoi
are identical by definition. If $r=1$ in
characteristic zero, the mixed geometric
topos is the arrow topos of B.4; its constant
cohomology is that of a point. If $r\geq2$
in characteristic zero, or $r\geq1$ in
positive characteristic, B.6 makes the mixed
comparison effective, so its cohomology is
that of $\mathbb P^1$. Its prime-to-$p$
degree-two group is $M^{(p')}(-1)$ and its
$p$-primary positive cohomology is zero.
This gives the table. The two-stage geometric
cohomology complex and Hochschild--Serre
give exactly the stated arithmetic cone,
including possible nontrivial transgression.
The one-puncture characteristic-zero case
now exhibits the final distinction: its
mixed-to-filled map is acyclic, while its
curve-to-mixed map misses $H^2$.
\end{proof}

\paragraph{C.6 (Proved): actions, Galois coverings, and the quotient square.}\label{C6}
\textit{Manuscript: LM \S16, folio 85--86, scans 172--174.}
For a discrete group $\Gamma$, coherent actions on
a Grothendieck topos $\mathcal E$ are equivalent,
as a two-category of data, to a topos $\mathcal X$
with a $\Gamma$-torsor $V$ and an equivalence
$\mathcal E\simeq\mathcal X_{/V}$.
One may take $\mathcal X=\mathcal E^\Gamma$, the
topos of equivariant objects. There is a
structure map $\mathcal E^\Gamma\to\beta\Gamma$,
and an equivariant morphism $\mathcal E\to\mathcal F$
is equivalently a morphism of these quotients
over $\beta\Gamma$. Its square with the torsor
covers is two-cartesian. If $\mathcal E=\beta C$
is multi-Galois, the quotient is
$\beta(C\rtimes\Gamma)$ for the Grothendieck
construction of the coherent action. It is
connected exactly when $\Gamma$ is transitive
on the nonempty set $\pi_0C$.

\begin{proof}
The underlying-object functor
$\mathcal E^\Gamma\to\mathcal E$ preserves finite
limits and all colimits. Equivariant objects
form a Grothendieck topos, for example by
equivariant sheaves on the disjoint union
of the translates of a small generating
site. The image of the regular torsor of
$\beta\Gamma$ is $V=\coprod_{\gamma\in\Gamma}1$
with its translation action. An equivariant
object over $V$ is recovered from its fibre
over the identity and the translated copies
of that fibre. This gives
$\mathcal E^\Gamma_{/V}\simeq\mathcal E$.
Conversely the effective descent of objects
along the epimorphism $V\to1$ identifies
$\mathcal X$ with the equivariant objects
of $\mathcal X_{/V}$; the torsor equation
$V\times V\simeq\Gamma\times V$ identifies
the descent cocycle with the action law.
The same descent applies to inverse image
functors and their two-cells. Pulling back
the same universal torsor gives the
two-cartesian square of scan~173.

For a groupoid action, an equivariant
presheaf on $C$ is exactly a presheaf on
its Grothendieck construction: arrows are
a group element followed by an arrow in
$C$, with the coherence of the action in
the composition formula. All these arrows
are invertible. Its connected components
are the $\Gamma$-orbits of $\pi_0C$,
proving the final assertions.
\end{proof}

\paragraph{C.7 (Refuted literally, corrected for profinite operators).}\label{C7}
\textit{Manuscript: LM \S16, folio 85, scans 172.}
For infinite profinite $\Gamma$, the regular torsor
of C.6 is not an object of
$\beta_{\mathrm{cont}}\Gamma$. The precise arithmetic
analogue uses continuous actions and a
\emph{protorsor}, or a torsor for the localic
group $\Gamma$. Precisely, in the connected profinite
Galois setting a continuous coherent action is encoded by
an exact sequence of profinite groups
\[
 1\longrightarrow G\longrightarrow E\longrightarrow\Gamma
 \longrightarrow1.
\]
Its quotient is $\beta_{\rm cont}E\to\beta_{\rm cont}\Gamma$;
base change along the canonical point
$\Set\to\beta_{\rm cont}\Gamma$ is $\beta_{\rm cont}G$.
Changing base points replaces this sequence by the
equivalent profinite groupoid extension. This definition,
componentwise with a continuous $\Gamma$-set of components,
is the continuous action meant in the arithmetic applications.
It has finite-quotient descent, but is not an ordinary slice
by a discrete regular $\Gamma$-set. No assertion about an
undefined abstract action of the underlying discrete group
on an arbitrary topos is implicit here.

\begin{proof}
The stabilizer of a point in the regular
action is $1$, which is not open in an
infinite profinite group. Thus the regular
set is not continuous. More decisively,
the canonical point $\Set\to
\beta_{\mathrm{cont}}\Gamma$ cannot be a
slice by an object giving the regular
cover: every nonempty transitive object
has open stabilizer, whereas this point
has trivial isotropy. For the corrected
statement, an $E$-set restricts to a $G$-set together with
transport over $\Gamma$, and the multiplication in $E$
is exactly its coherent action law. Conversely that
transport and its coherence recover the continuous
$E$-action. This proves the quotient assertion by descent.
For a precise base-change calculation, write $p:E\to\Gamma$.
For every open normal subgroup $N$ of $E$ there is an exact
sequence of finite groups
\[
 1\to G/(G\cap N)\to E/N\to\Gamma/p(N)\to1.
\]
The finite torsor calculation in C.6 identifies
$\beta(E/N)\mathbin{\times}_{\beta(\Gamma/p(N))}\Set$
with $\beta(G/(G\cap N))$. The subgroups $p(N)$ are cofinal
among the open normal subgroups of $\Gamma$: for $V$ open
normal, choose $N\subset p^{-1}(V)$. Similarly $G\cap N$
is cofinal among open normal subgroups of $G$. Take coherent
inverse limits of these finite action topoi, using
\citep[VI, \S8]{SGA4II}. Bicategorical limits commute with
this finite pseudopullback by their universal property,
giving $\beta_{\rm cont}E\times_{\beta_{\rm cont}\Gamma}\Set
\simeq\beta_{\rm cont}G$. The pro-cover is represented by
the compatible torsors $\Gamma/p(N)$, not by the regular
set $\Gamma$. For components, each orbit of a continuous
action of the compact group $\Gamma$ on a discrete set is
finite; apply the same calculation to its finite groupoid
and then take the disjoint sum over the orbits.
Equivalently a localic torsor has the regular compact
group as fibre. This is the qualification
needed for the parenthesis ``profini''
on scan~172.
\end{proof}

\paragraph{C.8 (Proved): the component factor and quotient compatibility.}\label{C8}
\textit{Manuscript: LM \S16, folio 87--88, scans 175--178.}
For $I=\pi_0C$ there is the canonical factorization
\[
 \beta(C\rtimes\Gamma)\longrightarrow
 \beta(I\rtimes\Gamma)\longrightarrow\beta\Gamma,
 \qquad \beta(I\rtimes\Gamma)\simeq(\beta\Gamma)_{/I}.
                                                               \tag{C.8.1}
\]
The first functor kills the fibre isotropy;
the second is faithful on groupoid arrows.
Over an orbit represented by $i$, the
intermediate topos is $\beta\Gamma_i$.
Equivariant quotient commutes with both
the mixed Artin gluing and the hole-filling
pushout, with their respective lax and
invertible commutation cells retained.
For profinite operators use continuous
$I$ and the interpretation of C.7.

\begin{proof}
The action groupoid $I\rtimes\Gamma$ has
arrows $\gamma:i\to\gamma i$. An equivariant
set over $I$ is precisely its presheaf,
giving the slice equivalence. Its arrows
map injectively to $\Gamma$. Its isotropy
at $i$ is $\Gamma_i$, and its components
are the orbits. The quotient functor
$C\rtimes\Gamma\to I\rtimes\Gamma$
is identity on the operator and forgets
the internal arrow in the fibre.
This is exactly the factorization of
scans~175--177. In the profinite case
continuity of $I$ makes each stabilizer
open, as required for the slice assertion.

An equivariant mixed object is an
equivariant $M$, an equivariant family
$N_s$, and an equivariant map $a_s$.
These are exactly the objects obtained
by gluing the quotient open and closed
topoi with the induced gluing functor.
The equivalence on maps is the same
componentwise description; it preserves
finite limits and all colimits. For
the filled diagram, a functor from the
quotient pushout to any topos is a
compatible cone from all the translates,
with its operator cocycle. Imposing the
pushout relations and imposing that
cocycle commute: both describe precisely
the same objects and natural
transformations. Equivalently, two
bicategorical colimits can be taken
in either order by their universal
property. This proves the assertion
in the slice over $\beta\Gamma$ and,
for discrete operators, in ordinary
topoi. Continuous descent gives the
profinite Galois version. For the
surjective hole-filling maps C.2's
closed-normal-closure proof also
applies directly to the quotient
extension groups. This proves the
claim of scan~178 with its two-category
and continuity hypotheses explicit.
\end{proof}

\subsection{Exterior homomorphisms and their transporter groupoids}

\paragraph{C.9 (Refuted, with replacement): the precise domain of exterior functoriality.}\label{C9}
\textit{Manuscript: LM \S12, folio 50--51, scans 103--106.}
For discrete groups, the Hom groupoid of
geometric morphisms $\beta G'\to\beta G$
has objects homomorphisms $u:G'\to G$ and
arrows $g:u\to v$ satisfying $v=c_g u$.
Its automorphism group at $u$ is
$C_G(u(G'))$. The same formula holds
for continuous homomorphisms of profinite
groups. Hence an exterior homomorphism
with trivial centralizer specifies a
geometric morphism up to unique
isomorphism, and inner changes of a
centre-free group model have unique
coherent identifications. However such
homomorphisms do not, without a further
restriction, form a category closed
under composition.

\begin{proof}
Pulling the regular torsor back in the
discrete case gives a $G$-torsor with
commuting $G'$-action. Choosing one
element identifies it with $G$ and
identifies that action with $u$.
Changing the element conjugates $u$;
equivariant maps of torsors give
exactly the displayed arrows. The
finite torsors give the same proof
in the profinite case. The stabilizer
formula follows immediately. An orbit
groupoid with trivial stabilizers is
contractible, so the exterior class
can be used without a distinguished
representative and all its comparison
cells are forced. This proves the
valid assertion on scans~103--105.

For the compositional qualification
take $N=S_3^3$, $H=N\rtimes S_3$ with
$S_3$ permuting the three factors, and
$K=S_3$. The centralizer of $N$ in
$H$ is trivial: a nontrivial permutation
changes a factor, and an element of
$N$ centralizing $N$ lies in its
trivial centre. The projection
$H\to K$ has centralizer $Z(K)=1$.
But the composite $N\to H\to K$ is
trivial and has centralizer $K$.
All three groups are centre-free.
The exhaustive finite calculation in
the program verifies orders $216$,
$1296$, and $6$, and the three
centralizer sizes $1,1,6$.

The strongest categorical replacement
is therefore either the full
bicategory of homomorphisms and
conjugators, which always composes,
or an explicitly composition-closed
subcategory in which every composite
still has trivial centralizer.
Surjections between centre-free
groups give one such subcategory.
On it the contractible representative
groupoids descend the classifying-topos
pseudofunctor to exterior data with
unique coherence. A graph of admissible
arrows alone does not justify that
conclusion for arbitrary composites.
\end{proof}

\paragraph{C.10 (Proved): the transporter construction and the square on scan 107.}\label{C10}
\textit{Manuscript: LM \S12, folio 51--52, scans 106--108.}
Let $\Lambda$ be the conjugacy class of
an injective $\kappa:T\to G$ satisfying
$C_G(\kappa(T))=\kappa(T)$. Define
$\Pi_\Lambda$ to have the homomorphisms
in $\Lambda$ as objects, and
\[
 \Hom(\kappa_1,\kappa_2)
   =\{g\in G: \kappa_2=c_g\kappa_1\}.
                                                               \tag{C.10.1}
\]
It is connected, its automorphism
groups are canonically the indicated
copies of $T$, and its forgetful
functor $\Pi_\Lambda\to BG$ sends an
arrow to $g$. This constructs
$B_\Lambda\to\beta G$ without choosing
a base loop.

For the convention actually written
on scan~107,
\[
 \Lambda=u\Lambda'\circ d,
                                                               \tag{C.10.2}
\]
and the self-centralizing hypotheses
there, the functor is
\[
 \kappa'\longmapsto u\kappa' d,
      \qquad g'\longmapsto u(g').
                                                               \tag{C.10.3}
\]
It gives the square of classifying
topoi together with its commutation
datum. The reversed arrow in the
sentence on scan~108 and the unprimed
$\kappa$ in the preceding formula are
type errors; the drawn square has
the direction in (C.10.3).

\begin{proof}
Composition in (C.10.1) is group
multiplication, so associativity and
inverses are immediate. The stabilizer
is the stipulated centralizer. Since
an abelian group is its own centralizer
inside itself, the hypothesis also
implies $T$ is abelian. The forgetful
functor is faithful, although in the
profinite case it is a slice induction
only when the relevant subgroup is
open, as Part A explains.

If $\kappa'_2=c_{g'}\kappa'_1$, then
$u\kappa'_2d=c_{u(g')}u\kappa'_1d$,
which proves (C.10.3) and its strict
commutation with the one-object
groupoid maps for a fixed $u$.
If $v=c_hu$, the components $h$
give a natural isomorphism between
these diagrams. The condition
$C_G(u(G'))=1$ makes $h$ unique.
For three representatives these
isomorphisms compose, since their
conjugators compose and are unique.
Thus descent over the exterior class
gives the stated square and its
two-cell. Taking disjoint unions over
the labels $i'$ gives the family
version, with the correctly typed
formula $\Lambda_{\tau(i')}=u\Lambda'_{i'}d_{i'}$.
This proves the construction rather
than treating an exterior homomorphism
as a chosen actual homomorphism.
\end{proof}

\paragraph{C.11 (Refuted, with replacement): degrees and uniqueness in scans 107--109.}\label{C11}
\textit{Manuscript: LM \S12, folio 52--53, scans 107--109.}
For any abelian group $T$, the injectivity and
self-centralizing hypotheses in (C.10.2) force the
integer scalar $d:T\to T$ to be an automorphism.
For $T=\mathbb Z$ or $\widehat{\mathbb Z}$ this forces
$d=\pm1$. The induced map on the
identified inertia groups is $d^{-1}$,
not $d$. In the prime-to-$p$ version
an integer $d$ may instead be a
unit of $\widehat{\mathbb Z}^{(p')}$.
The usual geometric convention for
a possibly ramified degree is
$u\Lambda'=\Lambda_{\tau(i')}d_{i'}$;
then the local maps and their
commutation cells are included as
data, as in B.12.

The two uniqueness assertions in
the margin are false without a
condition on normalizers and on
the indexing family. The exact
replacement is that, for each
given exterior composite, the
set of pairs $(i,d)$ satisfying
the displayed exterior equation
must have one element. A sufficient
intrinsic condition is that the
indexing family's allowed scalar
orbits are disjoint and each scalar
orbit map is injective. Fixing an
orientation excludes inversion
only when the normalizer respects
that orientation.

\begin{proof}
Write $a=u\kappa'$ and $\kappa=ad$.
Every element of $a(T)$ commutes
with $\kappa(T)$, so lies in
$\kappa(T)$. Thus $a=\kappa e$
for an endomorphism $e:T\to T$.
Since $\kappa$ is injective,
$\kappa=\kappa ed$ gives $ed=1$.
Integer scalars commute with every endomorphism of an
abelian group, so also $de=1$; this proves the assertion
for every abelian $T$ (continuously as well when $\kappa$
is a topological embedding). For the stated procyclic groups
their endomorphism rings are
$\mathbb Z$ or the corresponding
completion. Therefore $d$ is a
unit and $e=d^{-1}$. Formula
(C.10.3) on automorphisms is
exactly $e$. This shows why a
nonunit ramification degree cannot
be inserted into (C.10.2) under
its written hypotheses.

For degree nonuniqueness let
$G=D_\infty=\langle r,s\mid
s^2=1,\ srs^{-1}=r^{-1}\rangle$,
$T=\mathbb Z$, and $\kappa(1)=r$.
The centre of $G$ is trivial and
$C_G(\kappa(T))=\langle r\rangle$.
With $u=\id$, the exterior class
of $\kappa$ equals that of
$\kappa\circ(-1)$, so $d=1$ and
$d=-1$ both satisfy the same
equation. The action of integers
on $T$ is faithful, as demanded
on scan~109. For label
nonuniqueness take $G=F(a,b)$
and the two distinct conjugacy
classes $1\mapsto a$ and
$1\mapsto a^{-1}$. They are
self-centralizing and not
conjugate, by cyclically reduced
word normal form. The same
source loop and $u=\id$ satisfy
the equation with the first
label and $d=1$, or the second
label and $d=-1$. Thus an
injective label map alone does
not determine $\tau$. Finally
the asserted exact replacement
is simply the fibre, over $[u\Lambda']$, of the map
$(i,d)\mapsto\Lambda_i d^{-1}$ for (C.10.2),
or of $(i,d)\mapsto\Lambda_i d$ for the geometric convention:
disjointness and injectivity
are precisely what make that
fibre a singleton. This proves
the replacement and both
counterexamples completely.
\end{proof}

\paragraph{C.12 (Proved): the operator construction declined on scan 109.}\label{C12}
\textit{Manuscript: LM \S12, folio 51--53, scans 105--109.}
Let $\Sigma$ act externally on a
centre-free group $G$, on the
loop labels $I$, and on $T$ by
$\chi:\Sigma\to\Aut(T)$, with
$\sigma\Lambda_i=\Lambda_{\sigma i}\chi_\sigma$
as exterior data. Define
\[
 E=\Aut(G)\times_{\Out(G)}\Sigma.
                                                               \tag{C.12.1}
\]
Then $1\to G\to E\to\Sigma\to1$
is the canonical extension.
There is a groupoid over $B\Sigma$
whose objects are $(i,\kappa_i)$
and whose arrows are $e\in E$
over $\sigma$ satisfying
\[
 j=\sigma i,\qquad
 \kappa_j=e\kappa_i\chi_\sigma^{-1}.
                                                               \tag{C.12.2}
\]
Here $e$ acts through its
automorphism of $G$. Its fibre
is $\coprod_i\Pi_{\Lambda_i}$;
forgetting the loop gives a
functor to $BE$ over $B\Sigma$.
Its classifying topoi are the
equivariant square required in
\S12. Compatible exterior maps
with trivial centralizer, and
the loop and degree data of
C.10, induce maps of these
squares with all coherence
cells. Composition is understood
in C.9's precise sense.

\begin{proof}
Since $G$ has trivial centre,
$\Inn(G)=G$, so the kernel of
(C.12.1) is exactly $G$ and
surjectivity follows from the
definition of $\Out(G)$. The
formula (C.12.2) is stable under
composition: for $e$ over
$\sigma$ and $e'$ over $\tau$,
the resulting formula uses
$e'e$ and
$\chi_{\tau\sigma}=\chi_\tau\chi_\sigma$.
It has identities and inverses.
External compatibility of the
loops supplies an arrow over
each allowed $\sigma$, after
adjusting its lift by an inner
automorphism. The fibre over
the identity is exactly the
transporter groupoid. The
operator action reconstructed
from this fibration is coherent
by its associative composition;
no arbitrary lift of each
operator is needed.

For clarity, we verify also
functoriality of the extension.
Let $u:G'\to G$ have trivial
centralizer and commute with
the external $\Sigma$-actions.
For $e'\in E'$ over $\sigma$,
there is an $e\in E$ over
$\sigma$ such that
\[
                  e\circ u=u\circ e'.
\]
External compatibility gives
existence after an inner
adjustment. Two choices differ
by an element centralizing
$u(G')$, so coincide. Uniqueness
then proves that $e'\mapsto e$
preserves multiplication and
restricts to $u$ on the kernels.
Send $(i',\kappa')$ to
$(\tau(i'),u\kappa'd_{i'})$
and an arrow $e'$ to $e$.
Equivariance of $\tau$, invariance
of the degree function, and
commutation with the $T$-action
give (C.12.2) for the images.
All comparison cells are the
unique conjugators already
constructed in C.10. This
proves the operator version
that the manuscript postpones.
For profinite groups take
continuous automorphisms,
continuous extensions and
finite-quotient descent; an
external action is required
to be continuous in precisely
that extension sense. An
abstract homomorphism into an
un-topologized outer automorphism
group would not suffice.
\end{proof}
